% Part: first-order-logic
% Chapter: models-theories
% Section: introduction

\documentclass[../../../include/open-logic-section]{subfiles}

\begin{document}

\olfileid{fol}{mat}{int}
\olsection{مقدمة}

\begin{explain}
يُعَد تطور المنهج البديهي إنجازًا مهمًا في تاريخ العلم، وله أهمية
خاصة في تاريخ الرياضيات. وينطوي العرض البديهي لمجال ما على توضيح
أسئلة كثيرة: ما موضوع المجال؟ وما مفاهيمه الأساسية؟ وكيف ترتبط هذه
المفاهيم بعضها ببعض؟ وهل يمكن تعريف جميع مفاهيم المجال بدلالة هذه
المفاهيم الأساسية؟ وما القوانين التي تخضع لها هذه المفاهيم، والتي
يجب أن تخضع لها؟

لقد خُلق المنهج البديهي والمنطق أحدهما للآخر. فالمنطق الصوري يوفر
أدوات لصياغة النظريات البديهية، ولإثبات المبرهنات من بديهيات النظرية
على نحو مضبوط بدقة، ولدراسة خصائص جميع الأنظمة التي تحقق البديهيات
دراسة منهجية.
\end{explain}

\begin{defn}
تكون مجموعة من ال!!{sentence}s~$\Gamma$ \emph{مغلقة} إذا وفقط إذا
كان، متى $\Gamma \Entails !A$، فإن $!A \in \Gamma$. و\emph{غلاف}
مجموعة ال!!{sentence}s~$\Gamma$ هو
$\Setabs{!A}{\Gamma \Entails !A}$.

ونقول إن~$\Gamma$ \emph{مبنية بديهيًا من} مجموعة من الجمل~$\Delta$
إذا كانت $\Gamma$ غلاف~$\Delta$.
\end{defn}

\begin{explain}
يمكننا النظر إلى النظرية البديهية بوصفها مجموعة الجمل المبنية
بديهيًا من مجموعة بديهياتها~$\Delta$. وبعبارة أخرى، حين تكون لدينا لغة رتبة
أولى تحتوي على رموز غير منطقية للمفاهيم الأولية في العلم المبني
بديهيًا الذي نريد دراسته، مع مجموعة من ال!!{sentence}s تعبر عن
قوانين ذلك العلم الأساسية، أمكننا النظر إلى النظرية على أنها ممثلة
بجميع ال!!{sentence}s في هذه اللغة التي تستلزمها البديهيات. ويتراوح
ذلك من أمثلة بسيطة ليس فيها إلا مفهوم أولي واحد وبديهيات بسيطة،
كنظرية الترتيبات الجزئية، إلى نظريات مركبة كالميكانيكا النيوتنية.

والحقائق المنطقية المهمة التي تكسب هذا النهج الصوري إلى المنهج
البديهي أهميته هي الآتية. افترض أن $\Gamma$ نظام بديهيات لنظرية،
أي مجموعة من الجمل.
\begin{enumerate}
\item يمكننا أن نصرح بدقة متى يلتقط نظام بديهيات فئة مقصودة من
  ال!!{structure}s. فإذا كنا نهتم بفئة معينة من ال!!{structure}s،
  فإن نظام البديهيات~$\Gamma$ يلتقط تلك الفئة بنجاح إذا وفقط إذا
  كانت ال!!{structure}s هي بالضبط البنى~$\Struct M$ التي تحقق
  $\Sat{M}{\Gamma}$.
\item قد نفشل في هذا لأن ثمة~$\Struct M$ تحقق
  $\Sat{M}{\Gamma}$، لكن $\Struct M$ ليست من ال!!{structure}s
  التي نقصدها. وقد يدفعنا ذلك إلى إضافة بديهيات لا تصدق
  في~$\Struct M$.
\item إذا نجحنا، على الأقل، في جعل~$\Gamma$ صادقة في جميع
  ال!!{structure}s المقصودة، فإن الجملة~$!A$ تصدق في جميع
  ال!!{structure}s المقصودة متى كان $\Gamma \Entails !A$. وهكذا
  يمكننا استخدام أدوات منطقية (مثل طرائق ال!!{derivation}) لإظهار
  أن جملًا تصدق في جميع ال!!{structure}s المقصودة بمجرد إظهار أن
  البديهيات تستلزمها.
\item لا تكون في أذهاننا أحيانًا !!{structure}s مقصودة، بل نبدأ من
  البديهيات نفسها: فنبدأ ببعض المفاهيم الأولية التي نريدها أن تحقق
  قوانين معينة نرمز إليها في نظام بديهيات. ومن الأمور التي نريد
  التحقق منها فورًا ألا تتناقض البديهيات بعضها مع بعض؛ فإن تناقضت
  لم يمكن أن توجد مفاهيم تخضع لهذه القوانين، وكنا قد حاولنا بناء
  نظرية غير متماسكة. ويمكننا التحقق من عدم وقوع ذلك بإيجاد نموذج
  لـ~$\Gamma$. وإذا كانت لنظريتنا نماذج، أمكننا استخدام الطرائق
  المنطقية لدراستها، واستخدام الطرائق المنطقية أيضًا لبناء نماذج.
\item استقلال البديهيات سؤال مهم كذلك. فقد تكون إحدى البديهيات
  في الواقع نتيجة للبديهيات الأخرى، فتكون زائدة. ويمكننا إثبات أن
  بديهية~$!A$ في~$\Gamma$ زائدة بإثبات
  $\Gamma \setminus \{!A\} \Entails !A$. ويمكننا أيضًا إثبات أن
  بديهية ما ليست زائدة بإظهار أن
  $(\Gamma \setminus \{!A\}) \cup \{\lnot !A\}$ قابلة
  للإرضاء. وبهذه الطريقة، مثلًا، أُثبت أن مسلمة التوازي مستقلة عن
  سائر بديهيات الهندسة.
\item ومن الأسئلة المهمة أيضًا قابلية تعريف المفاهيم في نظرية:
  فاختيار اللغة يعيّن مم تتألف نماذج النظرية. لكن ليس من اللازم أن
  يمثل كل وجه من أوجه النظرية منفصلًا في نماذجها. فمثلًا، يعيّن كل
  ترتيب~$\le$ ترتيبًا صارمًا مقابلًا~$<$---فمتى أعطي أحدهما أمكننا
  تعريف الآخر. ولذلك لا يلزم أن يحتوي نموذج لنظرية تتضمن ترتيبًا
  كهذا \emph{أيضًا} على الترتيب الصارم المقابل. فمتى يمكن، بوجه
  عام، تعريف علاقة بدلالة علاقات أخرى؟ ومتى يستحيل تعريف علاقة
  بدلالة غيرها (ومن ثم يجب أن نضيفها إلى المفاهيم الأولية للغة)؟
\end{enumerate}
\end{explain}


\end{document}
